% =================================================================
% CONFIGURAÇÕES DO TEMPLATE - VARIANTE TESE
% Configuração completa para teses de doutorado e mestrado.
% Como usar:
%  - Preencha os placeholders abaixo (título, autor, etc.)
%  - Adicione/remova pacotes conforme a necessidade do seu trabalho
%  - Mantenha overrides aqui, sem editar config/sbc-template.sty
% Características desta variante:
%  - Suporte completo a capítulos e partes (usar com documentclass report/book)
%  - Índice, lista de figuras, lista de tabelas, lista de abreviaturas
%  - Glossários e índice remissivo (habilitados por padrão)
%  - Todos os pacotes disponíveis
%  - Suporte a banca examinadora e co-orientador
% =================================================================

% Pacotes essenciais
\usepackage{config/sbc-template}
\usepackage[T1]{fontenc}   % melhor hifenização/acentuação com pdflatex
\usepackage[utf8]{inputenc}
\usepackage{microtype}     % microtipografia (kerning/hifenização)
\usepackage[english,brazil]{babel} % idiomas: en + pt-BR
\usepackage{csquotes}      % citações tipográficas
\usepackage{graphicx,url}
\usepackage{hyperref}
\graphicspath{{assets/}}   % pasta padrão para imagens

% Pacotes adicionais (todos incluídos para teses)
\usepackage{listings}           % listagens de código
\usepackage{xcolor}             % cores (necessário para listings)
\usepackage{amsmath,amssymb}    % matemática
\usepackage{amsthm}             % teoremas e ambientes matemáticos
\usepackage{booktabs}           % tabelas profissionais (\toprule, etc.)
\usepackage{tabularx}           % tabelas largas com colunas X
\usepackage{array}              % colunas avançadas em tabelas
\usepackage{enumerate}          % listas enumeradas flexíveis
\usepackage{enumitem}           % controle avançado de listas

% Glossário e índice (habilitados por padrão para teses)
\usepackage[acronym,toc]{glossaries} % glossário com acrônimos
\makeglossaries
% Para usar: \newacronym{label}{abrev}{significado} e depois \gls{label}
% Compile com: pdflatex + makeglossaries + pdflatex

\usepackage{imakeidx}           % índice remissivo
\makeindex
% Para usar: \index{termo} e compile com: pdflatex + makeindex + pdflatex

% BibLaTeX (opção B2 - descomente se preferir ao invés de BibTeX):
% \usepackage[backend=biber,style=authoryear,doi=true,isbn=false,url=true]{biblatex}
% \addbibresource{bibliography/sbc-template.bib}

% Minted (opcional, para código com realce de sintaxe via Pygments):
% Requer: \usepackage{minted} e compilar com TEXFLAGS="-shell-escape".
% \usepackage{minted}

% Pacotes adicionais para teses
\usepackage{longtable}          % tabelas longas que quebram páginas
\usepackage{multirow}           % células mescladas em tabelas
\usepackage{subcaption}         % subfiguras e subtabelas
\usepackage{appendix}           % controle de apêndices

% Estilo de código (exemplo; ajuste ou remova)
\definecolor{codegray}{RGB}{248,248,248}
\lstdefinestyle{code}{
    backgroundcolor=\color{codegray},
    basicstyle=\footnotesize\ttfamily,
    numbers=left,
    numbersep=5pt,
    frame=single,
    breaklines=true,
    breakatwhitespace=true,
}

% Configurações de formatação para teses
% Espaçamento entre capítulos no sumário
\usepackage{tocloft}
\setlength{\cftbeforechapskip}{0.5em}
\setlength{\cftbeforesecskip}{0.2em}

% Numeração de equações por capítulo (descomente se necessário)
% \numberwithin{equation}{chapter}
% \numberwithin{figure}{chapter}
% \numberwithin{table}{chapter}

% Hyperlinks (preenche metadados do PDF)
\hypersetup{
    colorlinks=true,
    linkcolor=black,
    urlcolor=blue,
    citecolor=blue,
    pdftitle={PhD Thesis},
    pdfauthor={John Smith},
    pdfsubject={[Área de Pesquisa]},
    pdfkeywords={[palavras-chave, separadas, por vírgula]},
    bookmarks=true,
    bookmarksopen=true,
    bookmarksnumbered=true,
    pdfstartview=FitH,
}

% Comandos de placeholders (consumidos em \title, \author e \address)
\newcommand{\DocumentTitle}{PhD Thesis}
\newcommand{\DocumentSubtitle}{[Subtítulo (opcional)]}
\newcommand{\AuthorName}{John Smith}
\newcommand{\AuthorInstitution}{University Department}
\newcommand{\AuthorLocation}{Boston -- MA -- USA}
\newcommand{\AuthorEmail}{john@example.com}
\newcommand{\AdvisorName}{[Nome do Orientador]}
\newcommand{\AdvisorTitle}{[Prof. Dr./Dra.]}
\newcommand{\CoAdvisorName}{[Nome do Co-orientador (se houver)]}
\newcommand{\CoAdvisorTitle}{[Prof. Dr./Dra.]}
\newcommand{\ThesisType}{[Tese de Doutorado / Dissertação de Mestrado]}
\newcommand{\ThesisDate}{[Mês] de [Ano]}
\newcommand{\DefenseDate}{[Data da Defesa]}

% Metadados padrão (para projetos que usam este arquivo como entrada principal)
\title{\DocumentTitle}
\author{\AuthorName\inst{1}}
\address{\AuthorInstitution\\
  \AuthorLocation\\
  \email{\AuthorEmail}
}

% Configurações gerais
\sloppy
\setlength{\parskip}{6pt}
\setlength{\parindent}{1.5cm}  % parágrafo com recuo (padrão para teses)

% Definições de ambientes para teoremas (descomente se necessário)
% \theoremstyle{plain}
% \newtheorem{theorem}{Teorema}[chapter]
% \newtheorem{lemma}[theorem]{Lema}
% \newtheorem{proposition}[theorem]{Proposição}
% \newtheorem{corollary}[theorem]{Corolário}
%
% \theoremstyle{definition}
% \newtheorem{definition}{Definição}[chapter]
% \newtheorem{example}{Exemplo}[chapter]
%
% \theoremstyle{remark}
% \newtheorem{remark}{Observação}[chapter]
